\section{Proofs}\label{app:proofs}

\paragraph{\Cref{prop:contraction}} By the triangle inequality,
$\|\frac1K\sum_k r_{k,c}\|\le\frac1K\sum_k\|r_{k,c}\|=\|r_c\|+\frac1K\big(\sum_{k\in H_c}\Delta_k+\sum_{k\notin
H_c}\Delta_k\big)$, and the two sums are $m_c\bar G_c$ and $-(K-m_c)\bar D_c$. Equality holds if and only if all vectors
are non-negative multiples of one another. \hfill$\square$

\paragraph{\Cref{lem:adam}} A non-holder's batches contain no sample of class $c$, so its data loss does not depend on
$r_c$ and the gradient of the regularised loss is $\lambda r_c$. Adam's update $\eta\hat m/(\sqrt{\hat v}+\epsilon)$ is
unchanged when the whole gradient sequence is multiplied by $a>0$, up to the $\epsilon$ term, so the step does not
depend on $\lambda$ while $\lambda|r_{c,j}|\gg\epsilon$. A client keeps its Adam state across rounds while its
parameters are reset to the global ones, so a single step can exceed $\eta$. \hfill$\square$

\paragraph{The step ratio (\cref{eq:kappa}), a heuristic} Per coordinate an Adam step is $\eta\psi$ with
$\psi=\hat m/\sqrt{\hat v}$. For a gradient of constant sign $\psi$ tends to that sign; in general
$|\psi|\le(1-\beta_1)/\sqrt{1-\beta_2}$, about 3.2 for the default $\beta_1,\beta_2$~\cite{kingma2015adam}. A non-holder
sees a decay gradient of constant sign, so $\psi\approx\mathrm{sign}(r_{c,j})$ on each of its $q_k$ steps. Assume that
a holder moves the row only on the $q_kf_{kc}$ steps whose batch contains class $c$, with $|\psi|\le1$, and neglect
its other steps. Under flat averaging the drift per round is
$\Delta r_{c,j}=\frac{\eta}{K}\big[\psi_h\sum_{k\in H_c}q_kf_{kc}-\mathrm{sign}(r_{c,j})\sum_{k\notin H_c}q_k\big]$,
which points toward zero for every coordinate and every holder direction when $\kappa_c<1$. The approximation neglects
momentum carried between sparse holder steps and the drift of $\hat v$. \hfill$\square$

\paragraph{\Cref{prop:reparam}} (1)~$f+a_k=Wu(h)+(b+a_k)$, and $b\mapsto b+a_k$ is a translation with identity
Jacobian, so the gradients agree at corresponding points; induct over the steps. (2)~$\alpha f_c=\tilde w_c^\top
u+\tilde b_c$ with $\tilde w_c=\alpha w_c$, $\tilde b_c=\alpha b_c$; the gradient in $\tilde w_c$ is $1/\alpha$ times
that in $w_c$, so a step of size $\eta$ on $w_c$ is a step of $\eta\alpha^2$ on $\tilde w_c$, and the body sees the
same logits. \hfill$\square$

\paragraph{\Cref{prop:invariance}} Addition is associative and commutative. In particular the class-balanced linear
discriminant of \cref{eq:er} is computed exactly by one round of secure aggregation. \hfill$\square$

\paragraph{\Cref{prop:variance}} $\frac1C\sum_c(\LSE(f)-f_c)=\LSE(f)-\bar f$ and
$\KL(u\|\mathrm{softmax}f)=-\log C+\LSE(f)-\bar f$; expanding $\LSE$ around constant logits gives the quadratic term.
\hfill$\square$

\begin{lemma}[alternation, quadratic losses]\label{prop:alternation}
Let the synthetic and real losses be quadratic with Hessians $H_s,H_r\succeq0$ and minimisers $\theta_s,\theta_r$, and
let $P_s=(I-\eta H_s)^{S_s}$, $P_r=(I-\eta H_r)^{S_r}$ for $S_s$ synthetic and $S_r$ real steps per round. Let $0\le\eta\lambda\le1$ for every eigenvalue $\lambda$ of $H_s$ and $H_r$. If
$\|P_rP_s\|<1$, the rounds converge to $\theta^\star=\theta_r+\Omega(\theta_s-\theta_r)$ with
$\Omega=(I-P_rP_s)^{-1}P_r(I-P_s)$. For commuting Hessians the synthetic weight along eigendirection $i$ is
$\omega_i=\rho_{r,i}(1-\rho_{s,i})/(1-\rho_{r,i}\rho_{s,i})\le\rho_{r,i}$, which equals one along directions that the real loss
leaves free and the synthetic loss does not.
\end{lemma}
\emph{Proof.} One round is the affine map $\theta\mapsto P_rP_s\theta+P_r(I-P_s)\theta_s+(I-P_r)\theta_r$, a contraction
when $\|P_rP_s\|<1$; its fixed point has the stated form. With commuting Hessians all matrices are diagonal in a common
eigenbasis. \hfill$\square$

\begin{lemma}[what a mean anchor can express]\label{prop:w1}
Let the class-$c$ sampler be $X_c=\mathrm{ReLU}(\mu_c+g(Z,A\mu_c))$ with $Z$ shared across classes and $g$
$L$-Lipschitz in its second argument. Then $W_1(G_c,G_{c'})\le(1+L\|A\|_2)\|\mu_c-\mu_{c'}\|$.
\end{lemma}
\emph{Proof.} Couple both classes through the same $Z$; ReLU is 1-Lipschitz. \hfill$\square$

In units of the root-mean-square within-class spread, two class means are on average $\sqrt{2FC/(C-1)}$ apart (over
distinct pairs). When $F$ is small the anchor carries little class information, and the factor $1+L\|A\|_2$ is not
controlled; this is why the gate of \cref{sec:method-gen} falls back to class rows.

\begin{lemma}[least-displacement moment correction]\label{prop:mc}
Among all maps $T$ with $\E\,T(X)=\mu_c$ and $\E\|T(X)-\mu_c\|^2=V_c$, where $X$ has mean $\tilde\mu_c$ and spread $\tilde V_c$,
the map $T^\star(x)=\mu_c+\sqrt{V_c/\tilde V_c}\,(x-\tilde\mu_c)$ minimises $\E\|T(X)-X\|^2$. With exact covariances instead of
spreads the minimiser is the linear map attaining the Gelbrich bound~\cite{gelbrich1990}.
\end{lemma}
\emph{Proof.} With $Y=T(X)$, $\E\|Y-X\|^2=V_c+\tilde V_c+\|\mu_c-\tilde\mu_c\|^2-2\E\langle Y-\mu_c,X-\tilde\mu_c\rangle$, and the
Cauchy--Schwarz inequality bounds the last expectation by $\sqrt{V_c\tilde V_c}$, with equality if and only if
$Y-\mu_c=\sqrt{V_c/\tilde V_c}(X-\tilde\mu_c)$. The deployed map adds a final ReLU and clips the scale, so it matches the
moments only approximately. \hfill$\square$

\section{Experimental details}\label{app:protocol}

\paragraph{Models} The extractor, the ten head architectures (CNN-1 to CNN-10) and GeFL-F's conditional VAE follow
GeFL-F~\cite{kang2025gefl}; heads are trained by SGD with learning rate 0.1. The extractor is one convolutional block
followed by $2\times2$ max pooling: one $3\times3$ convolution with three output channels on\whennum{main/mnist/FedECS-LT/bal}{ MNIST,}\whennum{main/fmnist/FedECS-LT/bal}{
Fashion-MNIST,}\whennum{main/svhn/FedECS-LT/bal}{ SVHN,} OrganAMNIST, EMNIST and FER-2013, and two convolutions with three and ten output
channels on CIFAR-10, CIFAR-100, DermaMNIST, BloodMNIST\whennum{nat/gtsrb/FedECS-LT/bal}{, GTSRB} and PathMNIST, each
convolution followed by batch normalisation and ReLU. Rounds of the three stages ($T_{\mathrm{I}}$/$T_{\mathrm{II}}$/$T_{\mathrm{III}}$):\whennum{main/mnist/FedECS-LT/bal}{
20/100/50 on MNIST,}\whennum{main/fmnist/FedECS-LT/bal}{ 20/100/50 on Fashion-MNIST,}\whennum{main/svhn/FedECS-LT/bal}{ 50/100/100 on SVHN,} 50/200/100 on CIFAR-10 and CIFAR-100; with natural imbalance,
20/100/50 on OrganAMNIST and EMNIST and 50/100/100 on DermaMNIST, BloodMNIST, PathMNIST\whennum{nat/gtsrb/FedECS-LT/bal}{,
GTSRB} and FER-2013.

\paragraph{Reproduction check} Before the comparison, the Flower implementation was checked against our earlier,
single-process implementation (which follows GeFL's code) on MNIST-LT and CIFAR-10-LT with seed 3. The client splits
were identical (every client's class counts), every client training step matched weight for weight in unit tests, and
the balanced accuracy of GeFL-F and \ours was within the larger of two points and three seed standard deviations of
the earlier runs. The earlier CIFAR-10 reference of \ours was an older configuration, so its tolerance was three
points.

\paragraph{Head architectures} CNN-1 to CNN-10 have convolutional blocks of widths (16), (16, 32), (20, 40), (10, 20),
(16, 32, 64), (20, 40, 80), (10, 20, 40), (16, 32, 64, 128), (20, 40, 80, 100) and (10, 20, 40, 80); the widths do
not increase monotonically from CNN-1 to CNN-10. Each block is a $3\times3$ convolution, batch normalisation (the
naturally imbalanced datasets\whennum{main/fmnist/FedECS-LT/bal}{, Fashion-MNIST}\whennum{main/svhn/FedECS-LT/bal}{, SVHN} only), ReLU and $2\times2$ max pooling, followed by
a linear last layer.

\paragraph{FedDF} After each Stage-III round the server distils every architecture for 20 SGD steps (batch 128, learning
rate 0.05, no momentum) on fresh generated features with uniformly drawn labels; the target is the mean softmax of
the ten architectures, and batch-normalisation statistics are kept. FedDF's original method distils on unlabelled
public data instead.

\paragraph{GeFL's own setting} In GeFL's setting (balanced training data divided evenly among 10 clients,
\R{repro/cifar10/GeFL-F/seeds} seeds, our earlier implementation), GeFL-F reaches \R{repro/cifar10/GeFL-F/final}\% on
CIFAR-10 after the last round and \R{repro/cifar10/GeFL-F/best}\% with GeFL's measure, against
\R{repro/cifar10/published}\% read from GeFL's Fig.~4~\cite{kang2025gefl}, which reports each model's best round. An
earlier configuration of \ours (one synthetic epoch per round and an older head training) reaches
\R{repro/cifar10/FedECS-LT/final}\% and \R{repro/cifar10/FedECS-LT/best}\%. On CIFAR-100, which GeFL does not report,
GeFL-F reaches \R{repro/cifar100/GeFL-F/final}\% and that configuration \R{repro/cifar100/FedECS-LT/final}\%.

\paragraph{Evaluation} Models are evaluated every five rounds on the test set before and after aggregation; the final
model is the one after the last round (no selection on the test or the validation set). Methods with personal
heads (LG-FedAvg, FedGH, FedGen and the other HtFLlib methods) are scored with each client's head after its last local
training; HtFLlib scores LG-FedAvg, FedGH and FedGen right after the server's update instead. FedProto and FedTGP classify by the distance to
class prototypes; their probabilities for the calibration metrics are a softmax of the negative mean squared distances.

\paragraph{Cost} Communication is the bytes of the model and statistics messages per client, with a secure sum
counted as eight bytes per summed value and client, without the protocol's key and share messages and without the
test evaluation; the readout's sums are included. Time is the wall-clock time of a run on a T4 GPU, which in most
notebooks ran a second run at the same time. Computation is the forward floating-point operations per image, with
training counted as three forward passes.

\paragraph{Software and hardware} \ours and every method of the main setting run in Flower
1.42~\cite{beutel2020flower}, a simulation with the Ray backend, one simulation per run. All runs used Kaggle notebooks
with two NVIDIA T4 GPUs, in most notebooks with two runs per GPU at a time.

\paragraph{HtFLlib ports} Every client keeps the frozen Stage-I extractor and starts from its Stage-I head body,
followed by a parameter-free pooling to 512 dimensions and a linear layer, as in HtFLlib's models. Training uses
HtFLlib's run-script settings (SGD, learning rate 0.01, batch 10) with 5 local epochs per round for the rounds of
Stage III; personal models persist across rounds. FedGen's label counts are the federation's class counts.

\section{Further results}\label{app:more}

Per-dataset tables of every metric are generated from the same runs (\cref{tab:app-cifar10,tab:app-cifar100}).
Besides the metrics of \cref{sec:setup-metrics} they give macro-F1, macro precision, the multi-class Matthews
correlation coefficient~\cite{gorodkin2004mcc}, the geometric mean of the class recalls (each floored at 1\%, computed
for each model and then averaged) and macro AUC-PR; the class-wise form of the calibration
error~\cite{kull2019classwise} is logged as well. The figures of this appendix show further views of the same runs
(\cref{sec:results} cites each).

\begin{widetable}[t]
\centering\small
\caption{CIFAR-10-LT (imbalance ratio 100, 10 clients, Dirichlet $\alpha=0.5$): main metrics (\%); \ours with its readout, its heads only, other methods as published;
mean $\pm$ standard deviation over seeds 3--5. Tail: 40\% least frequent classes; G-mean with recalls floored at 1\%.}
\label{tab:app-cifar10}
\slottab{main_cifar10_main}{\begin{tabular}{lcccccccc}\toprule
Method & Bal. acc. & Macro-F1 & Precision & MCC & G-mean & Tail rec. & Worst rec. & AUC-PR\\\midrule
\phrow{FedECS-LT}{8}\phrow{GeFL-F}{8}\phrow{[every other method]}{8}\bottomrule\end{tabular}}
\end{widetable}

\begin{widetable}[t]
\centering\small
\caption{CIFAR-100-LT: main metrics (\%), as \cref{tab:app-cifar10}.}
\label{tab:app-cifar100}
\slottab{main_cifar100_main}{\begin{tabular}{lcccccccc}\toprule
Method & Bal. acc. & Macro-F1 & Precision & MCC & G-mean & Tail rec. & Worst rec. & AUC-PR\\\midrule
\phrow{FedECS-LT}{8}\phrow{GeFL-F}{8}\phrow{[every other method]}{8}\bottomrule\end{tabular}}
\end{widetable}

\whenfig{fig_forest}{%
\begin{figure}[t]
\centering
\slotfig{fig_forest}{col}{Fig.~39 of the CIFAR-10-LT report: difference of each method to FedECS-LT with 95\% CI.}
\caption{Balanced accuracy of \ours minus that of each method on CIFAR-10-LT (points; heads as trained for every
method, \ours included; mean of the paired
differences over seeds 3--5 with 95\% $t$ confidence interval; right of the dotted line \ours is better).
Imbalance ratio 100, 10 clients, Dirichlet $\alpha=0.5$.}
\label{fig:forest}
\end{figure}}

\whenfig{fig_recall_size}{%
\begin{figure}[t]
\centering
\slotfig{fig_recall_size}{col}{Fig.~10 of the CIFAR-100-LT report: per-class recall against the class's training
images.}
\caption{Recall of every class against its number of training images in the federation (CIFAR-100-LT, imbalance
ratio 100, 10 clients, Dirichlet $\alpha=0.5$; heads as trained, before the readout of \ours; each point is one class, mean over seeds 3--5).}
\label{fig:recall-size}
\end{figure}}

\whenfig{fig_hmt}{%
\begin{figure}[t]
\centering
\slotfig{fig_hmt}{col}{Fig.~14 of the CIFAR-10-LT report: head / medium / tail recall per method.}
\caption{Recall of the head (30\% most frequent), medium (the next 30\%) and tail (40\% least frequent) classes, per
method (CIFAR-10-LT, imbalance ratio 100, 10 clients, Dirichlet $\alpha=0.5$; heads as trained, before the readout of
\ours; mean $\pm$ standard deviation over seeds 3--5).}
\label{fig:hmt}
\end{figure}}

\whenfig{fig_shared_rounds}{%
\begin{figure}[t]
\centering
\slotfig{fig_shared_rounds}{col}{Fig.~71 of the CIFAR-100-LT (ratio 100) one-shared-model report: balanced accuracy
over the rounds.}
\caption{Test balanced accuracy over the communication rounds with one shared ResNet-8 (CIFAR-100-LT, imbalance ratio
100, 20 clients, 40\% per round, Dirichlet $\alpha=0.5$; mean over seeds 3--5). For \ours the rounds are those of the extractor and then of
the head.}
\label{fig:shared-rounds}
\end{figure}}

\whenfig{fig_realgen}{%
\begin{figure}[t]
\centering
\slotfig{fig_realgen}{col}{Fig.~22 of the CIFAR-10-LT report: t-SNE of real and generated validation features of
FedECS-LT.}
\caption{Real validation features (circles) and features generated by \ours (crosses) on CIFAR-10-LT, embedded
together by t-SNE (seed 3; same settings and caveats as \cref{fig:tsne}).}
\label{fig:realgen}
\end{figure}}

\whentab{fidelity}{%
\subsection{Fidelity of the generated features}\label{app:fidelity}

After Stage I the server trains a referee classifier (CNN-5 head, Adam, 40 epochs) on the features of its validation
images, which no client holds, and measures for each class the share of 300 generated features of that class that the
referee assigns to it. \Cref{tab:fidelity} gives the mean over the classes of each group. Chance is 10\% on
CIFAR-10-LT and 1\% on CIFAR-100-LT. The referee's accuracy on held-out real features was not logged, so the values
compare generators with each other, not with real data.

\begin{table}[t]
\centering\small
\caption{Generator fidelity (\%): the referee's per-class accuracy on generated features, mean over the classes of each
group (head: 30\% most frequent, tail: 40\% least frequent); mean $\pm$ standard deviation over seeds 3--5. FedDF uses
GeFL-F's generator, trained in its own runs.}
\label{tab:fidelity}
\begin{adjustbox}{max width=\linewidth}
\slottab{fidelity}{\begin{tabular}{llcccc}\toprule
Dataset & Generator & All & Head & Medium & Tail\\\midrule
\phrow{CIFAR-10}{5}\bottomrule\end{tabular}}
\end{adjustbox}
\end{table}}

\begin{widetable}[t]
\centering\small
\caption{GeFL's measure on the long-tailed benchmarks: each client model's best test accuracy (\%) at GeFL's
evaluation rounds of Stage III (the first, every tenth and the last), averaged over the ten models, as in GeFL's
code; mean $\pm$ standard deviation over seeds 3--5. The main tables report the last round instead. For \ours
only its heads: the readout is computed once, after the last round. FedDF: its client heads before the server's
distillation step, since the per-round logs hold no distilled heads; its value can therefore be below its last-round
value in \cref{tab:main}.}
\label{tab:app-gefl}
\slottab{main_gefl}{\begin{tabular}{lccccc}\toprule
Method & MNIST & Fashion-MNIST & SVHN & CIFAR-10 & CIFAR-100\\\midrule
\phrow{FedECS-LT}{5}\phrow{FedECS-LT (heads only)}{5}\phrow{GeFL-F}{5}\phrow{[every other method]}{5}\bottomrule\end{tabular}}
\end{widetable}

\begin{widetable}[t]
\centering\small
\caption{Accuracy (\%) on each client's own class mix (the test accuracy re-weighted to the class frequencies of the
model's own clients) on the long-tailed benchmarks; \ours with its readout, its heads only, other methods as
published; mean $\pm$ standard deviation over seeds 3--5.}
\label{tab:app-local}
\slottab{main_local}{\begin{tabular}{lccccc}\toprule
Method & MNIST & Fashion-MNIST & SVHN & CIFAR-10 & CIFAR-100\\\midrule
\phrow{FedECS-LT}{5}\phrow{FedECS-LT (heads only)}{5}\phrow{GeFL-F}{5}\phrow{[every other method]}{5}\bottomrule\end{tabular}}
\end{widetable}

\begin{widetable}[t]
\centering\small
\caption{Calibration on the long-tailed benchmarks: expected calibration error (\%, 15 bins; lower is better); \ours with its readout, its
heads only, other methods as published; mean $\pm$ standard deviation over seeds 3--5.}
\label{tab:app-calib}
\slottab{main_ece}{\begin{tabular}{lccccc}\toprule
Method & MNIST & Fashion-MNIST & SVHN & CIFAR-10 & CIFAR-100\\\midrule
\phrow{FedECS-LT}{5}\phrow{FedECS-LT (heads only)}{5}\phrow{GeFL-F}{5}\phrow{[every other method]}{5}\bottomrule\end{tabular}}
\end{widetable}

\begin{widetable}[t]
\centering\small
\caption{Fairness across clients on the long-tailed benchmarks: balanced accuracy (\%) of the weakest head
architecture; \ours with its readout, its heads only, other methods as published; mean $\pm$ standard deviation over seeds 3--5.}
\label{tab:app-worstarch}
\slottab{main_worst_arch}{\begin{tabular}{lccccc}\toprule
Method & MNIST & Fashion-MNIST & SVHN & CIFAR-10 & CIFAR-100\\\midrule
\phrow{FedECS-LT}{5}\phrow{FedECS-LT (heads only)}{5}\phrow{GeFL-F}{5}\phrow{[every other method]}{5}\bottomrule\end{tabular}}
\end{widetable}
